% concatenated from DST_.tex
\documentclass[a4paper,12pt]{amsart}
\usepackage[utf8]{inputenc}
\usepackage{mathtools}
%\usepackage{physics} %For divergence

\DeclareMathOperator{\RE}{Re}
\DeclareMathOperator{\IM}{Im}
\makeatletter
\@namedef{subjclassname@2020}{\textup{2020} Mathematics Subject Classification}
\makeatother

%\usepackage{cite}
\renewcommand{\baselinestretch}{1.1}
\allowdisplaybreaks[1]

\textwidth=460pt \evensidemargin=3pt \oddsidemargin=3pt
\marginparsep=8pt \marginparpush=8pt

\usepackage{amsmath,amsthm,amssymb}
\usepackage{hyperref}
\usepackage{amsfonts}
\usepackage{amsmath}
\usepackage{mathrsfs}
\usepackage{tabto}
\usepackage{changepage} %For adjustwidth
\usepackage{fullpage} %To decrease margins
\usepackage{graphicx}
%\usepackage{biblatex} %Imports biblatex package
\usepackage{enumitem} %For itemising
\newcommand{\ti}{\widetilde}
\newcommand{\la}{\lambda}
\newcommand{\La}{\Lambda}
\newcommand{\ga}{\gamma}
\newcommand{\ze}{\zeta}
\newcommand{\ity}{\infty}
\newcommand{\C}{\mathbb{C}}
\newcommand{\R}{\mathbb{R}}
\newcommand{\N}{\mathbb{N}}
\newcommand{\Z}{\mathbb{Z}}
\newcommand{\B}{\Big}
\newcommand{\F}{\mathfrak{F}}
\newcommand{\G}{\mathfrak{G}}
\newcommand{\A}{\mathfrak{A}}
\newcommand{\D}{\mathbb{D}}

\numberwithin{equation}{section}
%\numberwith{}
\newtheorem{theorem}{Theorem}[section]
\newtheorem{lemma}[theorem]{Lemma}
\newtheorem{corollary}[theorem]{Corollary}
\newtheorem{proposition}[theorem]{Proposition}
\newtheorem{revision}[theorem]{Revision}
%\theoremstyle{remark}
\newtheorem{remark}[theorem]{Remark}
\newtheorem{example}[theorem]{Example}
\newtheorem{definition}[theorem]{Definition}

\thanks {The research work of the first author is supported by ANRF(SERB) research grant TAR/2023/000197}

\begin{document}

\title[dynamics of bungee set of composite entire function]{results on dynamics of bungee set of composite entire functions in the eremenko-lyubich class}

\author[D. Kumar]{Dinesh Kumar}
\address{Department of Mathematics, Birla Institute of Technology Mesra, Ranchi, 
Jharkhand--835 215, India}
\email{dineshkumar@bitmesra.ac.in }

\author[S. Das]{Soumyajeet Das}
\address{Department of Physics, Birla Institute of Technology Mesra, Ranchi, 
Jharkhand--835 215, India}
\email{soumyajeetdas39@gmail.com }


%\author[D. Kumar]{Dinesh Kumar}
%\address{Department of Mathematics, Birla Institute of Technology Mesra, Ranchi, 
%Jharkhand--835 215, India}
%\email{dineshkumar@bitmesra.ac.in }

%\author[T. Nayak]{Tarakanta Nayak}
%\address{School of Physical Sciences, Department of Mathematics, IIT Bhubaneswar, Odisha--752050}
%\email{tnayak@iitbbs.ac.in}
%\maketitle	

 

\begin{abstract}
In this paper, we have discussed the dynamics of composite entire functions in terms of relationship between bungee set, escaping set and filled-in Julia set. We have established some relation between the dynamics of  composition of entire functions and the functions taken for composition. We have shown that the union of the bungee set of two entire functions contains the bungee set of the composite function. In addition, it is shown that the filled-in Julia  set of composite entire functions contains the filled-in Julia set of functions used for the composition. The results have been illustrated with several examples. We have mostly dealt with  permutable(commuting) functions.
% These results hold true when we have a composite function in which one function is a periodic translation of another.
\end{abstract}

\keywords{bungee set, wandering domain, Baker domain, repelling periodic point, rationally indifferent fixed point, escaping set, asymptotic value, singular set}

\subjclass[2020] {37F10,  30D05}


\maketitle


\section{Introduction}
When we iterate a holomorphic function $ f:\mathbb{C}  \rightarrow \mathbb{C} $, %or $ f: \mathbb{\hat{C}} \rightarrow \mathbb{\hat{C}} $
%($\mathbb{\hat{C}}=\mathbb{C}\cup{\{\infty\}} $
$n-$times, we observe that the points of the complex plane move differently, which we call the dynamics of the point. Here,  $f^n$  denotes the $n-$th iteration of  function  $f$, where $n \in \mathbb{N}$. As, not all  points have same dynamics, we therefore have different classifications of points of the complex plane $\mathbb{C}$; one classification is based on the nature of neighbourhood of a point $z_0$ (say); they are said to be in the Fatou set $F(f)$  (set of points whose neighbourhood  exhibits stable behaviour) and Julia set $J(f)$ (set of points whose neighbourhood exhibits chaotic behaviour). This was  broadly the classification we generally see when we talk about dynamics of points in complex plane. 

There is yet another classification of these points based on the nature of their orbits; viz.,  the filled-in  Julia set $K(f)$ (set of points whose orbits are bounded), escaping set $I(f)$ (set of points whose orbits tend to infinity), and bungee set $BU(f)$ (the orbits of such points neither remain bounded nor tend to infinity; only there exists a subsequence that remains bounded and another one that is unbounded). These sets  are disjoint among themselves and their union $BU(f)\cup I(f) \cup K(f)=\mathbb{C}$.  In \cite{osb2}, it was shown that if a Fatou component intersects with $BU(f)$ then that component becomes a wandering domain contained in $BU(f)$. 
%$K(f)$ comprises of  Fatou components in which orbits under iteration stays bounded, along with  repelling periodic points and rationally indifferent periodic points that  lies inside Julia set $J(f).$
It is to be noted that the bungee set and escaping set must have a non empty intersection with $J(f)$ (see \cite{e1, osb2}). The escaping set too has a non empty intersection with the Fatou set when the Fatou component is a Baker domain \cite{d1}, or when the Fatou set has a multiply connected wandering domain \cite{rs1}. Knowing all these sets for an entire function, we were curious about what happens when we compose two such functions and how these sets relate to the composition with those of individual functions. 
%We have seen this has been extensively used for rational maps. 
As composition of two  transcendental functions is again transcendental, we have tried to analyse what happens to the dynamics of transcendental functions under compositions.
% thus having all the sets we talked about a transcendental function in general. 
\\ 
Simplifying the situation further, we have taken transcendental functions $f ,g$  in the Eremenko-Lyubich class $\mathcal{B}$ which means that the functions of this class have a bounded set of singular values (i.e., the set of all  asymptotic values and critical values and their finite limit points is bounded). Recall the notion of an oscillatory wandering domain.
\begin{definition}
 Given a transcendental entire function $f$, a domain $U\subset F(f)$ is called an oscillatory wandering domain of $f$, if  $U\subseteq BU(f)$, that is, $\{\infty, d\}$ is a limit function of $\{f^n\}$ on $U$ for some $d\in \mathbb{C}$, \cite{marti2020wandering}.
\end{definition}
 Additionally, class $\mathcal{B}$ may have an intersection of the Fatou set with the bungee set, in other words there can be oscillating wandering domain which also lies inside the bungee set for $f\in\mathcal{B}.$ 
As our aim was to establish a relationship between different sets of functions and their compositions,  it would be very difficult to establish any such relation in the absence of some algebraic relation among these functions. Hence, we have focused on permutable (commuting) functions, that is, functions $f, g$ satisfying $f \circ g =g \circ f.$ Also, we have considered entire functions $f$ and $g$ related by $g = f+ C$ where $C$ is the period of $f$. 
The following results are proven in this paper:\\
Let $f,g \in \mathcal{B}$ be commuting functions. The bungee set of composed function  is contained in the union of the bungee set of functions used for composition. 
%The set of repelling periodic points and indifferent fixed points of the composed function denoted by $RP(f \circ g)$ contains the intersection of  repelling periodic points and indifferent fixed points of these functions; that is, $RP(f \circ g) \supset RP(f)\cap RP(g)$.\\ 
We know some results for exponential functions which belong to Speiser class $\mathcal{S}$ (such class of functions has finite number of elements in their singular set) \cite{sch}. Also, $\mathcal{S} \subset \mathcal{B}$. It was shown in \cite{baker2} that for $f= e^{\lambda z}, \la\neq 0$, a function $g$ that commutes  with $f$ is some finite iteration of $f$ itself, that is, $g=f^n$ for some $n\in\N$. Hence, 
%for such a case, all sets of composed functions are equal to the functions taken for composition.
in this case, the result clearly holds, in fact, we obtain $BU(f)=BU(g)=BU(f\circ g).$

\subsection{Observations}
While developing the relationships among the sets of composite functions and the individual functions, we also found
that the bungee set and escaping set exhibit similar characteristics for functions of class $\mathcal{B}$:
\begin{enumerate}
\item The closure of both these sets equals Julia set.
\item Both sets are neither open nor closed \eqref{Bug_nop_nclo}.
\item Both sets of composite function is contained in union of individual functions (to be proved later on).
\end{enumerate}

\section{Main results}
%-------------------------------------------------------------------------------------------------%
%Consider the Fatou map $f(z) = z+1+e^{-z}$\; and its linear translation $g(z)=f(z) +2\pi i.$
%Suppose $h(z)$ is a finite composition of $f$ and $g,$ i.e.
	 %$h(z)=f(z)^{l} \circ g(z)^{m}$ for some $l, m\in\N.$  \; This will lead to  the formation of a  %semigroup by the generators  $f$ and $g.$ \\
	%Now  $h(z)=f^{(m+l)} +2m\pi{}i,$ and
%$||f^{m+l}|^{n}-|2 n(m\pi) i|| \le |h^{n}(z)|.$
%\nonumber
%\\
%This shows when a point escapes to infinity for the sequence $\{f^n\},$ it also escapes to infinity under $\{h^n\}$ as well. It is known that the right half plane $U$ is contained in a Baker domain of $f$ cite{dinesh}, and so $U\subset I(f).$  Therefore, $U\subset I(h).$ \\
%Now $ \exists $ a point $z_{0}$ in the left half plane which is  pre-image of an escaping point in $U$ under $h,$ hence, $h(z)$ has points that escapes, are points from union of $U$ and  iterated pre-images of $U,$ i.e left half plane. This means escaping set of $I(h)$ has points from left half plane as well. For example, the image of the point $z_{0}=-1$  under $f(z)$ is $f(z_0)=e^{1}$ which under  iteration $\rightarrow \infty.$ \\
%So, this establishes that $U\subset$ I(G). \\
%However, if we take union of pre-images of the $U$(right half plane): {lim$_{n \rightarrow \infty }$ $\cup^{n}_{i=1} f^{-i}(U)$}, it almost fills the left half plane except uncountably many rays in the left half plane \cite{schleicher, dinesh}.
%\\
%We now introduce the notion of Bungee set of a semigroup $G$ to be denoted by $BU(G).$\\
%\begin{definition}
%$z_{0} \in BU(G)$ if for each sequence $h_{n}\in G,\exists$ two subsequences $\{m_{k}\},\,  \{n_{k}\}$ both tending to infinity, such that $|h_{m_k}| \rightarrow \infty \,\&\, |h_{n_k}|\leq K$ for some $K.$ It follows that $BU(G)$ is contained in the intersection of Bungee set of all the generators of the semigroup. 
%\end{definition}
%As $BU(G)$ has intersection definition, so, in general  $BU(G)$ can be the empty set.\\
%We now give some instances where  $BU(G)\neq \emptyset.$\\ 
%We may take general entire functions which are periodic translates to each other, i.e. suppose
 %$f(z)$ is an entire function and  $g(z)=f^{l}(z)+b$, where $b$ is the period of $f$. Then,
 %\begin{align*}
 %g^{i}&=f^{il}+b \\
 %g^{i}\circ f^{j}&=f^{il+j}+b\\
 %f^{i}\circ g^{j}&=f^{i+lj}\\  
 %\end{align*}
 %If $z_{0} \in BU(f)$ then,\\
 %\begin{align*}
 %h=g^{i}\circ f^{j}&=f^{il+j}+b\\
 %||f^{m_{k}(il+j)}|-|b|| \leq |h^{m_{k}}| \\
 %\end{align*}
%As $z_0\in BU(f),$ the left hand inequality tends to $\infty.$ Hence, $\{h^{m_k}\}$ tends to $\infty$ as well for subsequence $\{m_k\}$.
%Also,
%\begin{align*}
 %h=g^{i}\circ f^{j}&=f^{il+j}+b\\
 %|h^{n_{k}}|\leq |f^{n_{k}(il+j}|+|b|  \\
 %\end{align*}
 %Also, the right hand inequality is bounded. Hence, $\{h^{n_k}\}$ is bounded as well. This implies that $z_{0} \in BU(h)$.\\
 %Similarly, if $\psi$ is any other finite composition of elements of $G,$ then from the definition of Bungee set of semigroup,  $z_{0} \in BU(\psi)$ implies $BU(f)=BU(G)=BU(\psi)$. This example establishes that $BU(G)$ is a non empty set as $BU(f)\neq\emptyset$ by \cite{osb2}.
 %\\
 %We now provide an example of a transcendental semigroup $G$ for which the $BU(G)$ is an empty set, \cite{dk}.\\
 %We consider two families of functions in the Speiser class $\mathcal{S}$,
 %\\
 %$\mathcal{F}$ =\{$f_{\lambda,\xi}(z)=e^{-z+\lambda}+\xi :\lambda,\,\xi \in \mathbb{C}$,\,Re($\lambda) <0$ Re($\xi)\geq 1$\}\\
 %$\mathcal{F'}$ =\{$f_{\mu,\zeta}(z)=e^{z+\mu}+\zeta :\mu,\,\zeta \in \mathbb{C}$,\,Re($\mu) <0$ Re($\zeta)\leq -1$\}\\ It was shown that
%for each$f \in \mathcal{F}$, $J(f)$ is contained in the left half plane and 
%for each $g \in \mathcal{F'}$, $J(g)$ is contained in the right half plane.\\
%Since, these functions are of class $\mathcal{S}$, so they do not have wandering domains \cite{sch}.\\
%As a result, $BU(f) \subset J(f)$ and $BU(g) \subset J(g)$. And due to the nature of the two families, $J(f)\cap J(g) =\emptyset.$ So does $BU(f)\cap BU(g)$. Accordingly, by definition of $BU(G),$ it is an empty set.
 %\\
 %The above two examples suggests that when the entire functions $f$ and $g$ are connected by some kind of relation, then we get $BU(G) \neq \emptyset. $ \\
 %We now propose some questions and provide answers to them.
 %\\
%It was shown in \cite{osb2}, for each transcendental entire function $f, BU(f)\cap J(f)\neq\emptyset.$ We now ask similar question for a transcendental semigroup.\\
%\textbf{Question 1.} Can $BU(G)$ intersect $J(G)$?\\
%According to the definition of $J(G),\,$   $J(g) \subset J(G) \,\forall g \in G$ \cite{poon1}.  Also, we know $BU(g)\neq \phi$ and $BU(g)\cap J(g)\neq \phi$ \cite{osb2}.
%This implies that $BU(g)\cap J(g) \subset J(G).$ As $BU(G) \subset BU(g)$ $\forall$ $g\in G$, hence $BU(G)\cap J(G) \neq \phi$.   
 %\\
%Also, when the generators of the semigroup $G$ belong to class $\mathcal{S}$, then Fatou set of  $G,$ $F(G)$ does not have wandering domains \cite{dinesh}. 
 %From the definition of $J(G)$ we can also conclude that, $BU(G) \subset J(G)$ when the generators belonging to class $\mathcal{S}$. 
%\\ 
%---------------------------------------------------------------------------------------------------------------------
%
%Recall that a  component $U$ of $F(G)$ is called a \emph{wandering domain} of $G$ if the set $\{U_g:g\in G\}$ is infinite (where  $U_g$ is the component of $F(G)$ containing $g(U)$) \cite{dinesh1}.
%
%------------------------------------------------------------------------------------------------------------------
 %It was shown in \cite{osb2} that if $f$ is a transcendental entire function and $U\subset F(f)$ is any Fatou component. Then if $U\cap BU(f)\neq\emptyset,$ then $U$ is a wandering domain and $U\subset BU(f).$ We now ask similar question for a transcendental semigroup $G.$\\
%\textbf{Question 2.} What can we say about wandering domains of the semigroup $G$ if $U \cap BU(G) \neq \phi$ where $U\subset F(G)$ is a Fatou component.\\
%Let $U\subset F(G)$ be a  Fatou component satisfying $U\cap BU(G)\neq\emptyset.$ From \cite{osb2}, $U \cap BU(g)\neq\emptyset$ for all $g\in G.$ As a result $U\subset BU(g)$ for all $g\in G,$ and $U$ is contained in wandering domain of $F(g)$ for all $g\in G.$ This means that $U$ itself if a wandering domain of the semigroup $G$ and contained inside $BU(G).$ And since $BU(G) \subset \cap BU(g)\,\forall g\in G$, hence,  $\cap BU(g) \neq \phi$. 
% \\
%Hence, we have established that if $U\subset F(G)$ is a Fatou component and  $U \cap BU(G) \neq \phi,$ then $U\subset BU(G)$ and $U$ is a wandering domain of $G.$\\
%Let $z_0 \in U$ and since $z_0 \in F(G)$ also means,  $z_0 \in \cap F(g) \forall g \in G$. If $BU(g) \in F(g)$ we can also say that $F(G) \in \cup BU(g) \forall g \in G$, means $z_0 \in BU(g) \cup \forall g \in G$. Since it sits in its union means it also contains in $\cap BU( g) \forall g \in G $. So, $z_0 \in BU(G)$ So, Wandering domain $U \in BU(G)$.
%\\
%-------------------------------------------------------------------------------------------------%

\textbf{Property A} : $f$ (respectively $g$) has the property that it sends orbits escaping under $g$ (respectively, $f$) to $\ity.$\\

To begin with the commuting entire functions, we first understand the nature of these sets when functions satisfies \textbf{Property A}.

\begin{theorem} \label{Complete_invar_esp}
If $f$ and $g$ are commuting functions satisfying \textbf{Property A}, then $I(f)$ is completely invariant under $g$ and vice versa. 
\end{theorem}
\begin{proof}
Suppose $z_{0} \in I(f).$
Then, $|f^{n}(z_{0})| \rightarrow \infty$.
Now $|g(f^{n}(z_{0})|$ can either tend to $\infty$ or be bounded. Using  \textbf{Property A}, $|g(f^{n}(z_{0})| \rightarrow \infty $.
As $f$ and $g$ are commuting so, $|g(f^{n}(z_{0}))|$ = $|f^{n}(g(z_{0}))| \rightarrow \infty.$
This implies that  $I(f)$ is forward invariant under $g$.
It is easy to see that $g^{-1}(I(f))\subset I(f)$. 
Hence,  $I(f)$ is completely invariant under $g$. As $f$ and $g$ are commuting, their roles can be interchanged and hence, $I(g)$ is completely invariant under $f$ as well. 
\end{proof}

\begin{remark}
It follows from above result that   $I(f)$ and $I(g)$ are completely invariant under  $f \circ g$ as well.
\end{remark}
%\begin{lemma} \label{BU_no_asymp}
%Suppose $f$ and $g$ are commuting entire functions having no finite asymptotic values. Then $g(BU(f))\subset BU(f).$
% \end{lemma}
%
%\noindent 
%\begin{proof}
%Let $w\in BU(f).$ Then for subsequences $\{m_k\}$ and $\{n_k\}$ tending to $\ity$ we have 
%$|f^{m_k}(w)|< R$ and $|f^{n_k}(w)|\to\ity.$ It is clear that $|f^{m_k}(g(w))|< R_1$ for some $R_1$ and $|f^{m_k}(g(w))|$ tends to $\ity$ as $k\to\ity$ which implies that $g(w)\in BU(f)$ and hence the result.
%\end{proof}

\begin{theorem} \label{K_BU_no_asymp} 
Suppose $f$ and $g$ are commuting entire functions satisfying \textbf{Property A}. Then $K(f)$ and   $BU(f)$ are completely invariant under $g$ and vice versa.
\end{theorem}
\begin{proof}
$K(f)$ is in general forward invariant under $g,$ because if we take a point $w\in K(f)$, then we have $|f^{n}(w)|< R$ for some $R.$ It is clear that $|g(f ^{n}(w))|=|f ^{n}(g(w))|< R_1$ for some $R_1$  which implies that $g(w)\in K(f)$ that proves the result. Now, let $z_0 \in g^{-1}(K(f)).$ This implies that  $|f^{n}(g(z_0))| < L$ for some $L$ and because of commutation, $|g(f^{n}(z_0))|< L.$ Using  \textbf{Property A},  $|f^{n}(z_0)|<L_1$ for some $L_1$, which means $z_0 \in K(f)$. Thus, we obtain that $K(f)$ is completely invariant under $g$. Similarly, $K(g)$ is also completely invariant under $f$ as well.
\end{proof}

\begin{theorem} \label{K_BU_no_asymp1} 
Suppose $f$ and $g$ are commuting entire functions satisfying \textbf{Property A}. Then  $BU(f)$ is completely invariant under $g$ and vice versa.
\end{theorem}
\begin{proof}
Using the dynamical partition of the complex plane, we obtain $BU(f)= \mathbb{C} \backslash (I(f))\cup K(f)).$ We know from preceding argument that $K(f)$ is completely invariant under $g$ and from Theorem \ref{Complete_invar_esp}, $I(f)$ is completely invariant under $g.$ Hence, on combining these results, we  conclude that $BU(f)$ is also completely invariant under $g$. On similar lines, $BU(g)$ will  also be completely invariant under $f$. 
\end{proof}

\begin{remark}
It follows from above result that   $I(f),$ $I(g); K(f), K(g);$ and $BU(f), BU(g)$ are completely invariant under  $f \circ g$ as well.
\end{remark}

%\begin{lemma} \label{BU_no_asymp}
%Suppose $f$ and $g$ are commuting entire functions having no finite asymptotic values. Then $g(BU(f))\subset BU(f).$
% \end{lemma}
%

The next result provides a relation between the filled-in Julia set of composite entire functions and the filled-in Julia set of the individual functions.
\begin{lemma}\label{lem2.3}
The filled-in Julia set of the composite function $f\circ g$ contains the intersection of the filled-in Julia set of the individual functions which are commuting  and satisfying \textbf{Property A}.
\end{lemma}

\begin{proof}
Suppose that $z_0\in K(f)\cap K(g)$. Consider the array of  following sequences:\\
$\begin{array}{ccccc}
f(g(z_0)) &f^2(g(z_0))&f^3(g(z_0))\cdots\\
f(g^2(z_0))&f^2(g^2(z_0))&f^3(g^2(z_0))\cdots\\
\vdots&\vdots&\vdots\\
f(g^k(z_0))&f^2(g^k(z_0))&f^3(g^k(z_0))\cdots\\
\vdots&\vdots&\vdots

\end{array}$\\
Now, the first row of the array will converge to say, $l$ for, if it does not, then $g(z_0)\in I(f)$ which implies that $z_0\in I(f)$, a contradiction. Using continuity of $g$, we can say that second row will converge to $g(l)$ and so on. Let $X$ be the set containing all these limit points. The following are the possible cases for set $X$.\\
Case 1: $l$ is periodic point of $g$. In particular, suppose that $l$ is a fixed point of $g$. Hence, every sequence of the array will be bounded which shows that $z_0\in K(f\circ g)$.\\
Case 2: $X$ is infinite and bounded. Then, by Bolzano-Weierstrass theorem, $X$ has a limit point which also shows that $z_0\in K(f\circ g)$.\\
Case 3: $X$ is unbounded, so that there exists a subsequence which  escapes to infinity. This in turn implies that $z_0\in I(f\circ g)\subset I(f) \cup I(g)$ (see \cite{d1}) which leads to  a contradiction.\\
From all the above cases, we obtain that  $K(f)\cap K(g)\subset K(f\circ g).$
\end{proof}
%First we'll show that what kind of points will be Bounded even for commutation. So, 
%-------------------------------------------------------------------------------------------------%
%Let $z_0 \in K(f \circ g)$ i.e. $|f^n \circ g^n (z_0)|< K$ for some $K$. This implies $|g^n (z_0)|< K_1$ because if we take possibilities like $|g^n (z_0)| \rightarrow \ity$ and further on we apply $f^n$ to it, it shouldn't give a bounded number as these functions do not have finite asymptotic value. Similarly, if $g^n{z_0)}$ shows Bungee-behaviour, we can argue the following way as $f^n \circ g^n (z_0)=g^m \circ f^{2m} \circ g^{2m}(z_0)$ cannot be bounded since $f$ carries $BU(g)$ to itself when there are no finite asymptotic values. So, only points that'll be bounded would be some coming from $K(g)$. Similarly, these functions are commuting under commuting as well and this means $z_0 \in K(f)\cap K(g)$. This shows that if there's a point in $K(f \circ g)$, that must be a point from $K(f)\cap K(g)$ as well and therefore proving $K(f)\cap K(g) \subset K(f)\cap K(g)$.
%-------------------------------------------------------------------------------------------------%
%A simpler case will be to prove relation of Bungee set of composition of transcendental entire functions which are commuting under composition with those of Individual functions \textbf{What?}.
 
The next result connects bungee set of composite entire function with that of individual functions.
\begin{theorem}
Suppose $f$ and $g$ are commuting transcendental entire functions  satisfying \textbf{Property A}. Then, $BU(f \circ g)\subset BU(f)\cap BU(g).$
\end{theorem}

\begin{proof}
%We  establish that $BU(f \circ g)\subset BU(f)\cap BU(g).$

%Let $f$ and $g$ be two permutable entire functions such that both $f$ and $g$ do not have any finite asymptotic values. Then $BU(f\circ g)\subset BU(f)\cap BU(g)\subset BU(f)\cup BU(g)$.

 Suppose $z_0\in BU(f\circ g).$ Then,  there exists two subsequences $\{n_k\},\{m_k\}$ and a constant $M>0$ such that $(f\circ g)^{n_k}(z_0)\to\infty$ as $k\to\infty$ and $|(f\circ g)^{m_k}(z_0)|<M \mbox{ for all } k\in \mathbb{N}$. It is enough to show that $z_0\in BU(g)$ (as the proof of $z_0\in BU(f)$ follows on similar lines). In order to show  this, we prove that $z_0$ cannot be in $I(g)$ as well as $K(g).$ We divide the proof into several cases:\\
Case(i): Suppose that $g^{n_k}(z_0)\to\infty$ and $g^{m_k}(z_0)\to\infty$ as $k\to\infty$.\\
We can assume that $g^n(z_0)\to\infty$ as $n\to\infty$
(for, if there exists a subsequence $\{l_k\}$ such 
that $g^{l_k}(z_0)$ stays bounded as $k\to\infty$ then 
this will clearly imply that $z_0\in BU(g)$). Also, using Theorem \ref{Complete_invar_esp}, $f^k(I(g)\subset I(g)$ for every  $k\in\mathbb{N}$. In particular, $f^{m_k}(z_0)\in I(g)$ which implies that $(g\circ f)^{m_k}(z_0)\to\infty $ as $k\to\infty$ which leads to a  contradiction. \\
Case(ii): Suppose that $g^{n_k}(z_0)\to b$ and $g^{m_k}(z_0)\to\infty$ as $k\to\infty$. This again shows that $z_0\in BU(g)$ and the result follows.\\
Case(3): Suppose that $g^{n_k}(z_0)\to \alpha $ and $g^{m_k}(z_0)\to \beta$ as $k\to\infty.$   It follows that $z_0\in K(g)$ (for,  if there exists a subsequence $\{l_k\}$ for 
which $g^{l_k}(z_0)\to\infty$ as $k\to\infty$ then we are back to Case(2) 
 and hence $z_0\in BU(g)$). Using Theorem \ref{Complete_invar_esp},  $f(K(g)\subset K(g)$. In particular, $f^{n_k}(z_0)\in K(g)$, i.e., there exists $L>0$ such that $|g^n(f^{n_k}(z_0))|<L \mbox{ for all } n\in\mathbb{N}$. This implies that $|(f\circ g)^{n_k}(z_0)|<L \mbox{ for all } k\in\mathbb{N}$ which is again a contradiction. Thus, from the above cases, we conclude that $z_0\in BU(g).$  The proof of $z_0\in BU(f)$ follows on similar lines. 
Hence, we obtain that $BU(f\circ g)\subset BU(f)\cap BU(g)\subset BU(f)\cup BU(g)$.
\end{proof}
%\\
%Given the preceding proof, we ask a question:\\
%\textit{\textbf{Question:}}  What can we say about points of the set $BU(f) \cup BU(g)\setminus BU(f)\cap BU(g),$ where  will it be situated?\\
%\textit{\textbf{Answer:}} As the functions $f$ and $g$ do not have any finite asymptotic values, therefore,  points of the set $BU(f) \cup BU(g)\setminus BU(f)\cap BU(g)$ cannot intersect with $K(f \circ g).$ Also, $I(f \circ g) \subset I(f) \cup I(g)$ \cite{d1} and as $(BU(f)\cup BU(g)) \cap (I(f)\cup I(g))=\emptyset,$ it means that $BU(f) \cup BU(g)\setminus BU(f)\cap BU(g)$ cannot intersect with $I(f \circ g)$ either. This implies that the only place it will be contained is $BU(f \circ g)$. However, if it happens, then such points must be contained in $BU(f)\cap BU(g)$ implying that no such points exist. This is only possible if $BU(f)=BU(g)$. %when there are no finite asymptotic values for these commuting functions. 
%\\
Given the preceding proof, we ask a question:\\
\textit{\textbf{Question:}}  What can we say about points of the set $K(f) \cup K(g)\setminus K(f)\cap K(g),$ where  will it be situated?\\
%Similar question can be asked for filled-in Julia sets as well, (as it was shown in Lemma \eqref{lem2.3} that $K(f \circ g) \supset K(f) \cap K(g)$).
  Let $w_0 \in K(f) \cup K(g) \ \backslash \ K(f) \cap K(g)$. These points won't intersect with $I(f \circ g)$ (as $I(f \circ g) \subset I(f) \cup I(g)$ and $K(f)\cup K(g) \cap I(f) \cup I(g)=\emptyset$ \cite{d1}) nor with $BU(f \circ g)$ (as $BU(f \circ g)\subset BU(f)\cup BU(g)$ and $BU(f \circ g) \cap K(f) \cup K(g)= \emptyset)$ So, they will intersect with $K(f \circ g).$    
%\\
%Hence, given the following results:
%\begin{align} \label{equa_K_BU}
%K(f \circ g)&=K(f)=K(g) \nonumber \\
%&\text{and} \nonumber \\
%BU(f \circ g)&=BU(f)=BU(g)
%\end{align} 
%\noindent one can conclude that $I(f \circ g)=I(f)=I(g)$, for
% we know that $K(f)\cup I(f) \cup BU(f)=\mathbb{C}= K(f \circ g)\cup I(f \circ g) \cup BU(f \circ g)=K(g)\cup I(g) \cup BU(g)$ and from the previous results \eqref{equa_K_BU} we can conclude that $I(f \circ g)=I(f)=I(g)$.


\noindent The following examples illustrate the results established above.
%We now give an alternative argument to show this fact.\\
%Consider the array of following sequences:
%\begin{align*}
%&f(g(z_{0}))\tab &f(g^2(z_0))\tab &\cdots \tab &f(g^{n_k}(z_{0}))\cdots \\
%&f^{2}(g(z_{0}))\tab &f^2(g^2(z_0))\tab &\cdots  \tab &f^{2}(g^{n_k}(z_{0}))\cdots \\
%&. \tab &.\\
%&. \tab &.\\
%&. \tab &.\\
%&f^{l}(g(z_{0}))\tab &f^l(g^2(z_0))\tab &\cdots \tab &f^{l}(g^{n_k}(z_{0}))\cdots \\
%&. \tab &.\\
%&. \tab &.\\
%&. \tab &.\\
%&f^{n_k}(g(z_{0}))\tab &f^{n_k}(g^2(z_0))\tab &\cdots \tab &f^{n_k}(g^{n_k}(z_{0}))\cdots \\
%&. \tab &.\\
%&. \tab &.\\
%&. \tab &.\\
%\end{align*}
%Here all the rows approaches $\infty$ as $n_k \rightarrow \infty.$ As a result, the diagonal sequence also tends to $\infty.$ 
%Hence, we have proved that there is some intersection between $BU(f)$ and $BU(g),$ i.e. $BU(f)\cap BU(g)\neq \emptyset$.\\
%-------------------------------------------------------------------------------------------------%
%-------------------------------------------------------------------------------------------------%
%Now suppose, $w \in BU(f \circ g) \ \backslash \ BU(f)\cup BU(g)$. So, either 
%$w\in I(f) \cup I(g)$ or $w\in K(f)\cup K(g)$.
%\\
%For $I(f) \cup I(g)$:
%$|f^{m_k}\circ g^{m_k}(w)|< L $ for some $L$ so that $g^{m_k}(w)|< L_{1}$ for some $L_1.$ On similar lines, $f^{m_k}(w)<L_2$ for some $L_2.$ This is a contradiction and so $w \not \in I(f) \cup I(g)$.
%\\
%\\
%For $K(f) \cup K(g)$:
 %$|f^{n_k}\circ g^{n_k}(w)| \rightarrow \infty $ means $|f(g^{n_k}(w))| \rightarrow \infty $ also (since for the composition to tend to $\infty$, the rows should tend to $\infty$.) This implies, $|g^{n_k}(w)| \rightarrow \infty$ which is a contradiction. Hence,  in either case,  $BU(f \circ g$) will not be contained in $I(f) \cup I(g)$ or $K(f) \cup K(g)$. As a result, $BU(f \circ g\subset BU(f)\cup BU(g).$  \\
%-------------------------------------------------------------------------------------------------%
%--------------------------------------------------------------------------------------------------------------------
% Also let's take $'l' \in$ $BU(g)\cap BU(f\circ g)$. We know that $f(BU(g))$ is forward invariant because they are commuting type, so this proves that $'l'$ is also contained inside $BU(g)$ and no contradiction as such. Similarly for $f$, where k $\in BU(f)\cap BU(f\circ g) $. Now, let's take a point not in $BU(f)\cap BU(g)$ but in $BU(f)\ \& BU(f\circ g)$, so the point $z_0 \in BU(f)\cap BU(f\circ g)- (BU(f)\cap BU(g))$. So, $\exists$ a subsequence $m_k$ for which,
%$g^{m_k}\circ f^{m_k}(z_0) < R $, this means $f^{m_k}(z_{0}$ is also bounded as $m^{k} \rightarrow \infty$, and the subsequence  for which $g^{n_k}\circ f^{n_k}(z_0) \rightarrow \infty $ will also have $f^{n_k}(z_0) \rightarrow \infty $. Similarly for $g(z)$. So $BU(f\circ g) \subset BU(f) \cup BU(g)$.
 %-----------------------------------------------------------------------------------------------------
%\\
%\\
%\textbf{A Question cum paradox:}\\
% If we recall the arguments discussed above i.e if $z_0 \in BU(f \circ g)$ implies $z_0 \in BU(f)\cap BU(g)$ but then $BU(f \circ g)$ set can't have points which are from $I(f) \cup I(g)$. We also know that $I(f \circ g) \subset I(f) \cup I(g)$. Question is where will such points i.e. $BU(f) \cup BU(g)\ BU(f)\cap BU(g)$ will go?

\begin{example}
Consider $f(z)=1+z+e^{-z}$(Fatou map) and  $g(z)=1+z+e^{-z}+ 2\pi i.$ It can be easily checked that $f$ and $g$ satisfies \textbf{Property A}. 
%and hence $BU(f \circ g) \subset BU(f)\cup BU(g).$
\end{example}
%Now, $(f \circ g)^n(z_0)=f^{2n}(z_0)+2n\pi i$ and $g^{2n}=f^{2n}(z_0)+2n\pi i$. If $z_0 \in I(g)$ means $z_0 \in I(f \circ g)$. 
\noindent One may also consider the following functions for illustration of above result. They also satisfy \textbf{Property A}.\\
 $f_1(z)=1+z+e^{z}$  which would be commuting with $g_1(z)=1+z+e^{z}+ 2\pi i$;  $f_2(z)=z+\sin(z)$ which would  be commuting with $g_2(z)=z+ \sin(z)+2\pi.$
\\
% z+ cosh(z) which again commutes with its periodic translate.
%-------------------------------------------------------------------------------------------------%
%\begin{theorem}
%Commuting transcendental entire functions having no finite asymptotic value must have Bungee set equal.
%\end{theorem}
%-------------------------------------------------------------------------------------------------%
%-------------------------------------------------------------------------------------------------%
%We focus on points from $BU(f)\cup BU(g) \setminus BU(f)\cup BU(g)$. Because these functions do not have any finite asymptotic value, meaning such points cannot intersect with $K(f \circ g)$ also, $I(f \circ g) \subset I(f) \cup I(g)$ means that they cannot intersect with $I(f \circ g)$ either. This implies that the only place it will go to is $BU(f \circ g)$. However, if it happens, then such points must be from $BU(f \circ g)\subset BU(f)\cap BU(g)$ implying that no such points exist. This is only possible if $BU(f)=BU(g)$ when there are no finite asymptotic values for these commuting functions.\\ 
%-------------------------------------------------------------------------------------------------%
%-------------------------------------------------------------------------------------------------%
%Given this, we can say similar things for escaping sets of these functions.
%\begin{theorem}
%Commuting Transcendental entire functions having no finite asymptotic value must have Escaping sets equal.
%\end{theorem}  
%To prove this, we must first show that $I(f \circ g) \subset I(f) \cap I(g)$. Let us consider point $z_0 \in I(f \circ g)$, which means, $(f \circ g)^{n}(z_0)=f^{n}\circ g^{n}(z_0) \rightarrow \infty$. This implies that $g^{n}z_0$ either remains bounded or tends to infinity.\\ 
%Let us first examine the bounded case, where $g^{n}z_0$ is bounded means $z_0 \in RP(g)$ and the result from BAKER shows that $f(RP(g) \subset RP(g)$ so will $f^{n}(RP(g))$. Then, $g^{k_1}\circ f^m \circ g^{k_2}(z_0)$ where $k_1, k_2 , m \rightarrow \infty$ remains bounded (this can be said by the definition of $g(RP(g)$), which is a contradiction. Therefore, $f^n(z_0) \rightarrow \infty$. Similarly, we can say that for $f$, owing to the commutation property, $f \circ g$ can also be written as $g \circ f$. Therefore, point $z_0$ must belong to the intersection of these escaping sets. Hence, $I(f\circ g) \subset I(f) \cap I(g)$.
%-------------------------------------------------------------------------------------------------%

\section{Some results for functions of class $\mathcal{B}$}
Now, we will prove some results when the functions belong to  the Eremenko-Lyubich class $\mathcal{B}$. 
These functions do not have a Baker domain; therefore, the escaping set does not intersect with the Fatou set. In addition,  the bungee set of the function does not intersect with a wandering domain in which the iterates tend to $\infty.$ However, functions in this class may have oscillating wandering domain. 
%Therefore, the Julia set for functions of class $\mathcal{B}$ comprises of $RP(f), BU(f)$ and $I(f)$; that is, $J(f)=BU(f)\cup I(f)\cup RP(f)$. For commuting functions $f$ and $g$, we have $J(f \circ g)\supset J(f)\cup J(g)$ \cite{hua1, dinesh1}. However, when these are commuting functions of class $\mathcal{B}$ we have,  $J(f)=J(g)=J(f \circ g)$ \cite{hua1}. We also know that $g(J(f))\subset J(f)$ and $\overline{RP(f)}=J(f),\overline{I(f)}=J(f)$ \cite{bergweiler,EL}.

The next result gives a sufficient condition under which the bungee set is contained in the Julia set.
\begin{proposition}
The bungee set for functions of class $\mathcal{B}$ having no oscillatory wandering domain does not intersect with Fatou set, that is, $BU(f)\subset J(f).$ 
\end{proposition}
\begin{proof}
 For a function $f,$ we know that if $BU(f)\cap F(f)\not = \emptyset $, then such components are  oscillatory wandering domains \cite{osb2} and these components are contained inside $BU(f)$. As $F(f)$ has no oscillatory wandering domains, it means, $BU(f)$ does not intersect with $F(f)$. In other words, $BU(f)\subset J(f).$ 
\end{proof}
\begin{remark} \noindent This also means that $BU(f) \subset J(f)$  for functions not necessarily of $\mathcal{B}$ having no oscillatory wandering domain. Also, it is easy to see that $BU(f)$ is completely invariant by the complete invariance of $I(f)$ and $K(f)$. As a result, by the minimality property of Julia set, $\overline{BU(f)}=J(f)$. It is also known that for a function of $\mathcal{B}, J(f)=\overline{I(f)}$ \cite{EL}.
\end{remark}
\noindent The escaping set of a function is in general, neither open, nor closed \cite{e1}. We have similar kind of result for the bungee set.
\begin{proposition}\label{Bug_nop_nclo}
The bungee set for functions of class $\mathcal{B}$ having no oscillatory wandering domain is neither open nor closed. 
\end{proposition}
\begin{proof}
We can easily prove  the complete invariance of bungee set by noting that $BU(f)= \mathbb{C} \backslash (I(f) \cup K(f))$. Now, if we assume that bungee set is  closed then it will be a closed completely invariant set. This implies that $J(f) \subset BU(f)$ (because $J(f)$ is the smallest closed completely invariant set) which is a contradiction as $J(f)$ contains periodic points. Furthermore, if bungee set is open, it will imply that it has interior points which is also not possible (because $BU(f)\subset J(f)$ and bungee set having interior  means Julia set also has interior which is not possible for an entire function $f$ for which $F(f)\neq\emptyset$ \cite{beardon}). Hence, bungee set is neither open nor closed.
\end{proof}
The following result provides equality of bungee set, filled-in Julia set and escaping set of two functions in which one function is periodic translate of the other. Note that the two functions are non-commuting.
\begin{theorem}
Functions which are periodic translate of each other  have equal bungee sets, escaping sets and filled-in Julia sets.
\end{theorem}
\begin{proof}
We consider two functions $f$ and $g$ where $g$ satisfies $g(z)=f(z)+C.$  Here, $C$ is the period of $f$ (more generally, we can take $g=f^l(z)+C$, for some $l\in\N$). It can be easily seen that $g^{n}(z)=f^{n}(z)+C$. Again, it is easy to show that if a point is bounded under $f^n,$ it is also bounded under $g^n.$ Similarly, if a point  is escaping under $f^n,$ it is also escaping under $g^n$. Hence, $K(f)=K(g), I(f)=I(g)$ and  $BU(f)=BU(g)$.
\end{proof}
Such functions can be viewed as almost commuting type with a constant shift.

%Now moving forward, we will study functions that are  commuting  of class $\mathcal{B}$ and  having finite asymptotic values.
%-------------------------------------------------------------------------------------------------%
%\begin{theorem}
%Suppose $f$ and $g$ are permutable entire functions. Then $g(BU(f))\subset BU(f).$
%\end{theorem}
%Proof:\,
%Let $w\in BU(f).$ Then for subsequences $\{m_k\}$ and $\{n_k\}$ tending to $\ity$ we have 
%$|f^{m_k}(w)|< R$ and $|f^{n_k}(w)|\to\ity.$ It is clear that $|f^{m_k}(g(w))|< R_1$ for some $R_1$ %and $|f^{n_k}(g(w))|$ tends to $\ity$ as $k\to\ity$ which implies that $g(w)\in BU(f)$ and hence the result.  
%-------------------------------------------------------------------------------------------------%
%\begin{lemma} \label{BU_RP_BU}
%For commuting functions $f$ and $g$ of $\mathcal{B},$ we have $g(BU(f)) \subset BU(f)\cup RP(f).$ 
%\end{lemma}
%
%\begin{proof} 
%If we recall the definition of a bungee point, there must exist a subsequence which is bounded and another subsequence which is unbounded. Now it is obvious that when the subsequence which is bounded $f^{m}(z_0) \rightarrow l_0$ (say) and when acted upon by $g$ will also keep it bounded (since $f, g$ are entire functions). Let us assume that there is only one such bounded subsequence and $l$ unbounded subsequence which are as mentioned below:
%\begin{align*}
%&f^{m}(z_{0})\rightarrow l_0 \\
%&f^{n_1}(z_{0})\rightarrow \infty \\
%&f^{n_2}(z_{0})\rightarrow \infty \\
%%&f^{n_3}(z_{0})\rightarrow \infty \\
%&\vdots \\
%&f^{n_l}(z_{0})\rightarrow \infty
%\end{align*}
%Now, there are two possible cases when  $g$ acts on those unbounded subsequences.
%\\
%\textit{Case 1.} All these subsequence gets bounded, i.e.,
%\begin{align*}
%&g(f^{n_1}(z_{0}))\rightarrow w_1 = f^{n_1}(g(z_{0}))\rightarrow w_1 \\
%&g(f^{n_2}(z_{0}))\rightarrow w_2 = f^{n_2}(g(z_{0}))\rightarrow w_2\\
%%&g(f^{n_3}(z_{0}))\rightarrow w_3 = f^{n_3}(g(z_{0}))\rightarrow w_3 \\
%&\vdots \\
%&g(f^{n_l}(z_{0}))\rightarrow w_{l} = f^{n_l}(g(z_{0}))\rightarrow {w_l}
%\end{align*}
%If this happens, then $g(z_0)\in $  $K(f)=F(f)\cup RP(f)$. As $z_0 \in J(f)$  and we know that $g(J(f))\subset J(f)$ hence, $g(z_0) \in RP(f)$ for the above case. Now, let us understand what does this mean by $g(z_0) \in RP(f)$: It means that $f^p(g(z_0))=g(z_0)$ for some $p\in\N.$ This also means $g(f^p(z_0))=g(z_0)$. Because $f$ and $g$ are not invertible functions we have  $f^p(z_0)\not = z_0$. So, $g(z_0) \in RP(f) \ \& \ z_0 \in BU(f) $ can be possible.
%\\
%\textit{Case 2.} When at least one of the subsequence still remains unbounded. Let that be the $k-$th subsequence i.e., $g(f^{n_k}(z_0)) \rightarrow \infty$ i.e., $f^{n_k}(g(z_0))\rightarrow \infty $. $g(f^{m}(z_0))\rightarrow l_1,$ i.e., $f^{m}(g(z_0)) \rightarrow l_1 $. This shows that there are at least two subsequences and in this case $g(z_0)$ 
%belongs to $BU(f)$. 
%Therefore, $g(BU(f))$ is contained in either $BU(f)$ or $RP(f)$.
%\end{proof}
%
%\begin{lemma} \label{I_RP_I}
%$g(I(f)) \subset I(f) \cup RP(f)$ for commuting functions of $\mathcal{B}.$
%\end{lemma}
%\begin{proof}
%Suppose $z_0\in I(f)$ and let  $\lim_{n \rightarrow \infty} g(f^n (z_0)) \rightarrow l_0$ (say) i.e.,  $f^n(g(z_0)) \rightarrow l_0.$ This means $g(z_0) \in K(f).$ But $z_0 \not \in K(f)$ as it belongs to $I(f)$. So $g(z_0) \in RP(f)$ for the above case. Now, suppose $\lim_{n \rightarrow \infty} g(f^n (z_0)) \rightarrow \infty$ i.e., $f^n(g(z_0)) \rightarrow \infty$. Hence, $g$ sends $I(f)$ to either $RP(f)$ or $I(f)$. So, $g(I(f)) \subset I(f) \cup RP(f)$.
%\end{proof}
%
%\begin{proposition} \label{BU_compo_class_B}
%$BU(f \circ g) \subset BU(f) \cup BU(g)$ for commuting functions of class $\mathcal{B}$.
%\end{proposition}
%
%\begin{proof}
%Let $z_0 \in BU(f \circ g)\ \backslash (BU(f)\cup BU(g)) $. This implies that $z_0 \in K(f)\cup K(g)$ or $z_0 \in I(f) \cup I(g)$. By definition of bungee set, there are subsequences $\{n_k\}$  and  $\{m_k\}$ tending to $\infty$ satisfying
%\begin{align*}
%& |f^{n_k}(g^{n_k}(z_0))| \rightarrow \infty \\
%& |f^{m_k}(g^{m_k}(z_0))| \rightarrow l_0 \ \ \text{for some } l_0.
%\end{align*}
%First, we assume that $z_0\in K(g)$ which in turn implies $z_0 \in RP(g)$ (there is a clear contradiction in the $F(g)$ case because $F(g)=F(f \circ g)$), and as a result, $f^{n}(RP(g)) \subset RP(g)$  and $g^n(f^n(RP(g)))$ which will be bounded. As can be seen, there are no subsequences that are unbounded. Similarly, there are no such points in $K(f)$.  Therefore, $z_0 \not \in K(f) \cup K(g)$.\\
%Now, we assume $z_0 \in I(g)$ which means there are two possible cases
%\begin{align} 
%\label{case 1} & |f(g^{n}(z_0))| \rightarrow l \ \text{for some} \ l \\
%\label{case 2} & |f(g^{n}(z_0))| \rightarrow \infty  
%\end{align}
%where $n$ denotes an arbitrary index unrelated to $n_k$ and $m_k$. If we have case \eqref{case 2}, then from the array of sequences argument, we see that all subsequences of $(f\circ g)^{n}(z_0)$ tend to infinity. If we have case \eqref{case 1}, there are further possible cases: $l \in I(f)$ or $l\in K(f)$  or $l \in BU(f)$. Again, if we have the first two cases, then we observe that  entire  composite sequence $(f\circ g)^n$ is (un)bounded. The third case is also not possible. This is because using the previous Lemma \eqref{BU_RP_BU}, we have $|f^{n_k}(g^{n_k})(z_0)|$ or $|f^{m_k}(g^{m_k})(z_0)|$ which equals $|g^{n_k}(f^{n_k})(z_0)|$ or $|g^{m_k}(f^{m_k})(z_0)|$ will be either completely bounded or unbounded. Therefore, $z_0 \not \in I(g)$. On similar lines, $z_0\not\in I(f)$.  
%Therefore, $z_0 \not \in I(f) \cup I(g)$. Thus, we see that there are no points s.t., $z_0 \in BU(f \circ g)\ \backslash (BU(f)\cup BU(g))$.
%Let $z_0 \in BU(f \circ g)\ \backslash (BU(f)\cup BU(g)).$ This means $z_0 \in K(f)\cup K(g)$ or $z_0 \in I(f) \cup I(g)$. By definition of Bungee set, there are subsequences $\{n_k\}$  and  $\{m_k\}$ tending to $\infty$ satisfying
%\begin{align*}
%& |f^{n_k}(g^{n_k}(z_0))| \rightarrow \infty \\
%& |f^{m_k}(g^{m_k}(z_0))| \rightarrow l_0 \ \ \text{for some } l_0
%\end{align*}
%\baslineskip
%First, consider the subsequence $|f^{n_k} \circ g^{n_k}(z_0)| \rightarrow \infty.$ This implies that either $|g^n(z_0)|\rightarrow \infty$ or $w_0$ for some $w_0$ as $n_k \rightarrow \ity$. If we take the $w_0$ case first, %and $z_0 \in RP(f)\cup RP(g)
%then $z_0 \in K(g)=F(g) \cup RP(g)$ (say).
%(where $RP(g)$ represents set of repelling periodic points of the function $g$). 
%There is a clear contradiction in the $F(g)$ case because $F(g)=F(f \circ g)$. When $z_0 \in RP(g),$ also implies $f^{n}(RP(g)) \subset RP(g)$  and $g^n(f^n(RP(g)))$ are bounded. Hence, we cannot have subsequences tending to $\infty$ condition, which contradicts with $BU(f \circ g)$ definition itself. So, $z_0 \not \in K(g)$ and similar agrument can be given for points from $K(f)$.
%\\
%\\
%Let us now consider the case where $g^{n_k}(z_0) \rightarrow \infty$. This further raises the case of $f(g^{n_k}(z_0))$ either $f(g^{n_k}(z_0))$ will converge to $l$ (say) or will tend to $\infty$. Again, if it is an unbounded case, then by the array of sequences argument, it can be proven that all subsequences of composition $f^k \circ g^k(z_0) \rightarrow \infty$ as $k \rightarrow \infty$ which will again lead to a contradiction. Now, for the case where $f(g^{n_k}(z_0)) \rightarrow l$ (say), this means that $l \in BU(f)$, $K(f)$ or $I(f)$. If we have either cases $K(f)$ or $I(f)$, only (un)bounded sequences will occur upon iteration with $f^k ,k \rightarrow \infty$ which is again a contradiction. Let us take $l \in BU(f)$ case. This means 
%\begin{align*}
%& |f^{{m_k}-1}(l)| \rightarrow w \ \text{for some} \ w \\
%& |f^{{n_k}-1}(l)| \rightarrow \infty
%\end{align*}
%\nolinebreak
%therefore, both subsequences are possible. However, if we take $f$ inside $g^n(f(z_0)) \rightarrow l,$ it means $f(z_0) \in RP(g)$ (Not in $F(g)$ because we assume that the points are from $J(g)$). This also implies that $f^{k}(z_0) \in RP(g)$ where $k \rightarrow \infty$; And now $g^{k_1}(f^n(g^{k_2}(z_0))) \rightarrow w$ (say) where $k_1 + k_2 = n$ , $k_1 \rightarrow \infty$, $k_2 \rightarrow \infty$ and $n \rightarrow \infty$ and as a result we will not get a subsequence that  will tend to $\infty$ for the above case i.e, $f(g^n(z_0)) \rightarrow l,$ ($g^{k_1}(f^n(g^{k_2}(z_0))))$ sequence being bounded, also assures that  the limit function is in $RP(g),$ it does not catch an unbounded (sub)sequence. %Moreover, if $g(z_0) \in RP(f)$ be won't get any subsequence $\rightarrow \infty$.   Thus, $z_0$ does not exist. %but earlier when $f \& g$ had no asymptotic values we proved that $BU(f \circ g) \subset BU(f) \cup BU(g)$. Therefore, atleast the equality condition i.e $BU(f \circ g) = BU(f) \cup BU(g)$ must be fulfilled for functions of class $mathcal{B}$ having no asymptotic values and bounded critical values. But this doesn't prove that $BU(f \circ g) \supset BU(f) \cup BU(g)$ for functions of class $\mathcal{B}$ which are commuting. However, some points of I(f) or I(g) are  Bungee points for $f \circ g$.
%Hence, we have $BU(f \circ g) \subset BU(f) \cup BU(g).$  
%\end{proof}

%\begin{corollary}
%For commuting functions of class $\mathcal{B}, (I(f)\cup I(g)) \cap BU(f \circ g)=\emptyset$
%\end{corollary}
%\begin{corollary}
%There are no points $z_0 \in K(f)\cup K(g) \cap BU(f \circ g)$
%\end{corollary}
%\begin{proof}
%We can follow the similar lines of proof that is done in the preceeding Proposition \eqref{BU_compo_class_B}. Let $z_0 \in I(f)$ (say). Then $g(f^n(z_0)) \rightarrow \infty$ or $L$. If $g(f^n(z_0)) \rightarrow \infty,$ then by array of sequences argument, we conclude that   $g^n \circ f^n(z_0)\not \rightarrow M$ for some $M$ which means that we do not have a subsequence that remains bounded upon iteration of $f \circ g$. Moreover, if $g(f^n(z_0)) \rightarrow L$ then  $g(z_0)\in RP(f)$ and therefore, $f^n(g^n(f^n(z_0))) \rightarrow L_1$ (say). Thus, we never get  a subsequence that  tends to $\infty.$ Similarly, it can be argued when $z_0\in I(g)$.
%\end{proof}
%
%%Revised proof of $I(f \circ g) \subset I(f) \cup  I(g)$:}\\
%\begin{revision}
%Revised proof of $I(f \circ g) \subset I(f) \cup  I(g)$ for commuting functions of $\mathcal{B}.$ (considering the existence of bungee set)
%\end{revision}
%Let's consider two commuting functions of class $\mathcal{B}$ i.e, $f \circ g = g \circ f$.\\
%Functions of class $\mathcal{B}$ have the property that $J(f)=J(g)$ and therefore $J(f \circ g)=J(f)$ \cite{dinesh1}, which
%implies that  $\mathbb{C} \backslash J(f \circ g)= \mathbb{C} \backslash J(f)$\\
%$\implies F(f \circ g)= F(f)=F(g)$\\
%This implies that a point $z_0 \in J(f \circ g)$ can't belong to $F(g)$ (or $F(f)$), because that would imply that $z_0 \in F(f \circ g)$ as well, and that's a clear contradiction (since $I(f \circ g)$ is a component of $J(f)$). Also $g(J(f)) \subset J(g)= J(f \circ g)$. So, upon iteration it stays in $J(f\circ g)$.\\
%\begin{proof}
%The motivation for the proof comes from \cite{d1}. Let $z_0 \in I(f \circ g) \ \backslash (I(f)\cup I(g))$ i.e., $z_0 \in RP(f) \cup RP(g)$ or $z_0 \in BU(f) \cup BU(g)$. By definition of escaping set: $f^{n} \circ g^{n}(z_0)=(f \circ g)^n(z_0) \rightarrow \infty $. 
%\textbf{NOTE:} \textit{$g^n(z_0) \not \rightarrow \infty $ because $z_0 \not \in I(g)$}. 
%also there can be no such points s.t. $z_0 \in F(g)$, that means $z_0 \in F(f \circ g)$, a clear contradiction. \\
%Now, because functions of class $\mathcal{B}$, then we know $J(f)=J(g)$ and therefore $J(f \circ g)=J(f)$\\
%$\implies \mathbb{C} \backslash J(f \circ g)= \mathbb{C} \backslash J(f)$\\
%$\implies F(f \circ g)= F(f)=F(g)$\\
%Assume $z_0 \in BU(f)$, then $g(BU(f))$ goes to either $BU(f)$ or $RP(f)$ using Lemma \ref{BU_RP_BU}. 
%%where RP means repelling periodic point. 
%When it goes to $RP(f)$, we know $g(RP(f))$ will remain in $RP(f)$  and upon further iteration of $f \circ g$ will stay bounded i.e., $(g\circ f)^n(z_0) \rightarrow L$ as $n \rightarrow \infty$ which is  a contradiction. If it goes to $BU(f)$ then again, its the same case as discussed before when $g$ acts on it. In the end, it will either be an entirely bounded sequence or there would be an (un)bounded subsequence. The later case arises only when $g^{n}(BU(f)) \subset BU(f) \ \forall n \in \mathbb{N}$ (where $n$ is some arbitrary index). But by definition, escaping sets must have all the subsequences tending to infinity. Hence, we arrive at a contradiction. On similar lines, it can be shown that when $z_0\in BU(g)$ we again arrive at a contradiction. So, $z_0 \not \in BU(f)\cup BU(g)$.\\
%Now the case where $z_0 \in RP(f) \cup RP(g)$. Assuming $z_0 \in RP(f)$ and as discussed earlier  we will have all bounded subsequences of $f^{n} \circ g^{n}(z_0)$ which is a contradiction. Similar arguments can be given when $z_0\in RP(g).$ Therefore $z_0 \not \in RP(f)\cup RP(g).$ Hence, $z_0 \in I(f \circ g) \ \backslash (I(f)\cup I(g))$ is not possible. Thus, $I(f \circ g) \subset I(f)\cup I(g)$.
%means $z_0 \in K(g)$ i.e. $K(g)=F(g)\cup RP(g)$, 
 %(RP means repelling periodic point). 
 %Let $z_0 \in RP(g)$ (say). 
%By definition of periodic point we know that $g^m(z_0)=z_{0}$ where $m$ is finite. 
%Since $f(RP(g))$ goes to $RP(g)$, $f^{m}(RP(g))$ goes to $RP(g)$, and by definition $g(RP(g))$ goes to $RP(g)$.
%, so is $f^n(RP(g))$(But $z_0$ has to still be $\subset J(f)$ or $J(g)$, therefore it can't intersect with $F(g)$). 
%Now, let us consider $g^n \circ f^n( f^n \circ g^n (z_0)).$ This is same as   $g^n \circ f^n(RP(g))$ which is bounded, and therefore a contradiction, because such points $\not \rightarrow \infty.$ (Here we exploited the commutation relation of $f$ and $g$). 
%\end{proof}
%Now let's take the other case where $z_0 \in F(g)$, that means $z_0 \in F(f \circ g)$, its again a clear contradiction. 
%\noindent On similar lines, same argument can be given when we have $z_0 \in RP(f).$ There we can consider $f \circ g$ and its iteration. This means that no such $z_0$ exists which is in $I(f \circ g)$, and not in $I(f)\cup I(g)$. So, $I(f \circ g) \subset I(f)\cup I(g)$.
%Also Let's take a point $w_0 \in I(f)$(say) $g^{m} \circ f^{m}(w_0)$ where $m \rightarrow \infty$.So, either $g(f^m(w_0)) \rightarrow \infty$ or is bounded $l_1$(say) now this $l_1\in BU(g)$ or $I(g)$ or $K(g)$. and we can't say that where it will land and also this argument is similar for I(g). So, from all these evidences we can say that $I(f \circ g) \subset I(f) \cup I(g)$.  

%For a function $f,$ we know that if $BU(f)\cap F(f)\not = \emptyset $, then such components are  wandering domain \cite{osb2} and these components sit inside $BU(f)$. But in case of functions of class $\mathcal{B}$, these wandering domains $\not \rightarrow \infty $, so such wandering domains don't have any subsequence that $\rightarrow \infty.$ This means, $BU(f)$ doesn't intersect with $F(f)$ for class $\mathcal{B}$. This also means that $BU(f) \subset J(f)$ only for functions of class $\mathcal{B}.$ As a result, by the minimality property of Julia set, $\overline{BU(f)}=J(f)$.

%\\
% \textbf{Question: Will the repelling periodic of $f$ i.e. $RP(f)$ be a bounded set or unbounded}\\
% By definition we know that $\overline{RP(f)}=J(f)$. Now, we know that $J(f)$ is unbounded, this means that $RP(f)$ is also unbounded. This is because, closure means set with its limit points, if the limit points are finite means that that closure is finite, but if the limit points $\rightarrow \infty$ then only the closure set is unbounded. And also for the limit points are unbounded means the set itself is unbounded. So by this logic $RP(f)$  is unbounded. 
% \\

%\textbf{Proof for relation of $RP(f \circ g)$ with $RP(f)\ \& \ RP(g)$:}\\
%\begin{proposition}
%\textit{$RP(f \circ g) \supset RP(f) \cap RP(g)$ for commuting functions of class $\mathcal{B}$.}
%\end{proposition}
%\begin{proof}
%Let us consider a point $z_0 \in RP(f)$. Then, $f^n(z_0)$ is bounded and contained in $RP(f)$. Also, $g^n(RP(f)) \subset RP(f)$, therefore $f^n(g^n(f^n(z_0)))$ is also bounded as $n \rightarrow \infty .$ (Observe that, $z_0$ can't be in $ F(f \circ g)$  as $F(f \circ g)=F(f)$).  Hence, $z_0 \in RP(f \circ g)$ as well. On similar lines, if $w_0 \in RP(g)$, then $w_0 \in RP(f \circ g)$.\\
% Now the question is can there be points not in $RP(f) \cup RP(g)$ but in $RP(f \circ g)$? To answer that let us go through Lemma \ref{BU_RP_BU}  once again. Suppose $g(BU(f))$ goes to $RP(f)$ and upon further iteration with $(g \circ f)^n (z_0)$, this will be bounded means $z_0 \in RP(f \circ g)$. Similarly, there will be some points from $I(f)$ that will be bounded upon $g^n \circ f^n(z_0)$ as $n \rightarrow \infty$ as discussed in Lemma \ref{I_RP_I}. Again it has to be in $BU(g)$ and $I(g)$. Thus we conclude that we do have points that is in $RP(f \circ g)$ but not in  $RP(f) \cup RP(g)$.
%\end{proof}
%There's a simpler proof for this which goes as follows:
%\begin{proof}
%\begin{align}
%& BU(f \circ g) \cup I(f \circ g) \subset BU(f) \cup BU(g) \cup I(f) \cup I(g) \nonumber \\
%& \implies J(f \circ g) \backslash RP(f \circ g) = J(f) \backslash RP(f \circ g) \supset RP(f) \cup RP(g) \nonumber \\
%& \text{Taking complement on both sides} \nonumber \\ 
%& BU(f \circ g)^c \cap I(f \circ g)^{c} \supset BU(f)^c \cap BU(g)^c \cap I(f)^c \cap I(g)^c \nonumber\\
%& \implies RP(f \circ g) \supset RP(f) \cap RP(g)
%\end{align}
%\end{proof}

\section{Illustrations:}
We now illustrate  $BU(f \circ g)\subset BU(f)\cup BU(g)$ with some examples. Our first example is commuting rational functions. \\
Consider the functions $f(z)=z^2$ and $g(z)=1/z^2$. Observe that $f \circ g(z)=1/z^4$. $BU(f)= \emptyset$, $BU(g)=\{z: |z|<1 \cup |z|>1\}$ and $BU(f \circ g)=\{z:|z|<1 \cup |z|>1\}.$ Clearly, $BU(f \circ g)\subset BU(f)\cup BU(g)$.
% Also we know functions of  class $\mathcal{S}$ can be treated as rational maps so in that way also it makes sense to  have union-like definition.\\
\\
Our next example is motivated from  \cite{tw}. Let us take the case where $f= p(z)\cdot e^{h(z)}$ where $p(z) , \ h(z)$ are polynomials and $g(z)=a \cdot f(z)$ for some $a \in  \mathbb{C}\setminus \{0\}$. For simplicity, let us take
$f(z)=z\cdot e^{a^{(k+1)} \cdot z^k},\ k \in \mathbb{N}$ \&
$g(z)=a \cdot z \cdot e^{a^{(k+1)} \cdot z^k}$.\\
$f \circ g(z) =f(g(z))=g(z)\cdot e^{a^{(k+1)} \cdot g^k(z)}=aze^{a^{(k+1)}z^k}e^{a^{(2k+1)}z^ke^{ka^{(k+1)}z^k}}$\\
$g \circ f(z) =g(f(z))=a \cdot f(z)\cdot e^{a^k \cdot f^k(z)}= aze^{a^kz^k}e^{a^{2k}z^ke^{ka^kz^k}}.$ Now, for $f \circ g = g \circ f,$ we need to have $a^{2k+1}=a^k \implies a^k(a^{k+1}-1)=0,$ which means $a$ is $(k+1)$-root of unity except `1' itself because that's trivial. 
\\
Now, for the simplest case let us take: $k+1=2 \implies a =\pm 1$ Since $a=1$ is a trivial case so let us take $a=-1$, $f(z)=ze^{-z^2};\ g(z)=-ze^{-z^2}$. 
%Here as the orbit of any point under either of the function tends to $\infty, f(z) \rightarrow 0, g(z) \rightarrow 0 $
 Observe that  $z=0$ is an asymptotic value of $f$ and $g$. It is also  a fixed point of $f(z), \ g(z)$ since $f(0)=0$ and $g(0)=0$. Now for the multiplier $\lambda$, $f'(z)=e^{-z^2}(1-2z^2), f'(0)=1$ which is an indifferent fixed point. And $1^1=1$ means $0\in J(f)$ or $J(g)$ (since both of them are equal). As $\infty$ is a special case and belongs to $J(f)$ so, this proves  that a sequence which tends to  $\infty$ if it asymptotes to a finite asymptotic value which is a member of the Julia set of $f$. 
%Moreover, the functions which are of the form $f= p(z)e^{h(z)}$ have a finite asymptotic value  $'0'$ and it is in $J(f)$ (rationally indifferent fixed point or repelling periodic point).
%(rationally-indifferent fixed point). 
\\
Now some illustrations for the bungee set:\\
Let $f(z)=ze^{z^2};\ g(z)=-ze^{z^2}.$ This implies $f^2=ze^{z^2}e^{z^2e^{2z^2}} \ \& \ g^2=ze^{z^2}e^{z^2e^{2z^2}}.$
$f \circ g = -ze^{z^2}e^{z^2e^{2z^2}} = g \circ f$. And $(f\circ g)^2=g^4=f^4$(because as $(f \circ g)^2$ is done it does away with the negative term and rest terms are exactly same). So, $BU(f^4)=BU(g^4)=BU((f \circ g)^2)$ which implies $BU(f)=BU(g)=BU(f \circ g)$ (\textbf{NOTE:} If we have $f(z)=ze^{-z^2};\ g(z)=-ze^{-z^2}$, they also commute). We may also  take $f(z)=ze^{a^{(k+1)}z^k}$, $g(z)=aze^{a^{(k+1)}z^k}$ and $a^{(k+1)}=1$ for that matter). And this complies with the union relation with the composition.\\
Now, for the general case: Since $a$ is $k$-th root of unity, means $a^k=1$. So, $f(z)=ze^{a^k z^k}=ze^{z^k}$ and $g(z)=aze^{z^k}$. Now $g^k(z)=a^k  f^k(z)=f^k(z)$ implies $BU(f)=BU(g)$. Also, $(f \circ g)^{k}(z)=a^{k}f^{2k}(z)=f^{2k}(z)$. So, $BU(f \circ g)=BU(f)=BU(g)$. Analogously, if we take $f(z)=ze^{a^{(k+1)}z^k}$ and $\ g(z)=af(z)=aze^{a^{(k+1)}z^k}$, where again $a^{k+1}=1$, we can simplify it further as $f(z)=ze^{az^k}$ and $g(z)=aze^{az^k}$, now $g^2(z)=af(af(z)))=a^2f^2(z)$. By induction,  $g^{k+1}(z)=a^{(k+1)}f^{k+1}(z)$. Also, $f \circ g(z)= af^2(z)=g \circ f(z)$ and $(f \circ g)^2(z)=a^2f^4(z)$. Proceeding by induction, $(f \circ g)^{k+1}(z)=a^{(k+1)}f^{2(k+1)}(z)=f^{2(k+1)}(z)$. Here also, we can see that $BU(f\circ g)= BU(f)=BU(g)$.\\
Thus, functions which are of the form $f(z)=ze^{a^k z^k}$ where $a$ is $k$-th root of unity or $f(z)=ze^{a^{k+1} z^k}$, where $a$ is $(k+1)$-th root of unity and $g=af$  have the property that their bungee sets are same. Moreover, their escaping set and repelling periodic point sets are also the same. $J(f)=J(g)$ also holds true  because we are taking  $f= p(z)e^{g(z)}$ as discussed in \cite{tw}. 

%\section{Conclusion:}
%We have shown that entire functions with no finite asymptotic values which are also commuting 
%i.e., $f \circ g = g \circ f$ satisfies:\\
%$K(f)=K(g), \ I(f)=I(g), \ BU(f)=BU(g)$ and therefore, $K(f \circ g)=K(f), \ BU(f \circ g)=BU(f), \ I(f \circ g)=I(f)$. \\
%And for the functions of class $\mathcal{B}$ which are commuting \\
%%i.e $f \circ g = g \circ f$:\\ 
%$RP(f \circ g) \supset RP(f)\cap RP(g)$;\\
%$BU(f \circ g) \subset BU(f)\cup BU(g)$;\\
%$I(f \circ g) \subset I(f)\cup I(g)$.\\
%There are some other special cases where we have equality among these sets, like when $f$ is an exponential map. Also,  cases like $g=f+C$ where $C \in  \mathbb{C}$ is the period of $f$.
%Also, when functions are from class $\mathcal{B}$, $(BU(f)\cup BU(g)) \cap RP(f \circ g) \neq \emptyset $ because $g(BU(f)) \subset RP(f)\cup BU(f)$.
%Similarly,  $(I(f)\cup I(g)) \cap RP(f \circ g) \neq \emptyset$ because $g(I(f)) \subset RP(f)\cup I(f)$.
%Also, there are no points which are in $I(f) \cup I(g)$ as well as in $BU(f \circ g)$.

\begin{thebibliography}{00}
\bibitem{baker1} I. N. Baker, Limit functions and sets of non-normality in iteration theory, Ann. Acad. Sci. Fenn. Ser. A. I. Math. \textbf{467} (1970), 1-11. 
\bibitem{baker2} I. N. Baker, Wandering domains in the iteration of entire functions, Proc. London Math. Soc. \textbf{49} (1984), 563-576.
\bibitem{beardon} A. F. Beardon, \emph{Iteration of rational functions}, Springer Verlag, (1991).
\bibitem{bergweiler} W. Bergweiler, Iteration of meromorphic functions, Bull. Amer. Math. Soc. \textbf{29} (1993), 151-188.
\bibitem{e1} A. E. Eremenko, On the iteration of entire functions, Ergodic Theory and Dynamical Systems, Banach Center Publications \textbf{23}, Polish Scientific Publishers, Warsaw, (1989), 339-345.
\bibitem{EL} A. E. Eremenko and M. Yu. Lyubich, Dynamical properties of some classes of entire functions, Ann. Inst. Fourier, Grenoble, \textbf{42} (1992), 989-1020.
\bibitem{F2}  P. Fatou, Sur l'it\'eration des fonctions transcendantes Enti\`eres, Acta Math. {\bf 47} (1926), no.~4, 337-370. 
\bibitem{martin} A. Hinkkanen and G. J. Martin, The dynamics of semigroups of rational functions I, Proc. London Math. Soc. (3) \textbf{73} (1996), 358-384.

\bibitem{Hua} X. H. Hua, C. C. Yang, \emph{Dynamics of transcendental functions}, Gordon and Breach Science Pub. (1998).
\bibitem{hua1} X. Hua and X. Wang,  Dynamics of permutable transcendental entire functions, Acta Math. Vietnam. \textbf{27} (2002), 301-306.
\bibitem{dinesh1} D. Kumar and S. Kumar, The dynamics of semigroups of transcendental entire functions I,  Indian J. Pure Appl. Math. \textbf{46} (2015), 11-24.
\bibitem{dk} D. Kumar, On escaping sets of some families of entire functions and dynamics of composite entire functions, Math. Student, \textbf{84} (2015), No. 1-2, 87-94.
\bibitem{dinesh3} D. Kumar and S. Kumar, The dynamics of semigroups of transcendental entire functions II,  Indian J. Pure Appl. Math. \textbf{47} (2015), 409-423.
\bibitem{d1} R. Kaur and D. Kumar, Results on escaping set of an entire function and its composition,  Indian J. Pure Appl. Math. \textbf{52}(2021), 79-86.
\bibitem{tw} T. W. Ng, Permutable entire functions and their Julia sets, Math. Proc. Camb. Phil. Soc.\textbf{131} (2001), 129-138.
\bibitem{osb2} J. W. Osborne and D. J. Sixsmith, On the set where the iterates of an entire function are neither escaping nor bounded, Ann. Acad. Sci. Fenn. Math.  \textbf{41} (2016), 561-578.

\bibitem{marti2020wandering} D. M. Pete  and M. Shishikura, Wandering domains for entire functions of finite order in the Eremenko-Lyubich class, Proc. London Math. Soc. \textbf{120}  (2020), 155-191. 

\bibitem{poon1} K. K. Poon, Fatou-Julia theory on transcendental semigroups, Bull. Austral. Math. Soc. \textbf{58} (1998), 403-410.

\bibitem{rs1}  P. J. Rippon and G. M. Stallard, On questions of Fatou and Eremenko, Proc. Amer.
Math. Soc. \textbf{133} (2005), no. 4, 1119-1126.


\bibitem{sch} D. Schleicher, Dynamics of entire functions,  Lecture Notes in Mathematics  ((LNMCIME,volume 1998)) (2010), 295-339.

\bibitem{aps} A. P. Singh, On bungee sets of composite transcendental entire functions,  arXiv:2006.00208v1[math.CV] (2020) , pp. 1-7.
\end{thebibliography}
 

 \end{document}
 
